\section{Centralized Energy References}\label{sec:bounds}
We compute reference energies with full instance information and the same
power menu, slot duration, and service model. These references require
\emph{every packet} to meet its deadline. The subscript $0$ below denotes this
zero-loss requirement, distinct from the $\epsilon=0.05$ allowance for learned
policies.

\subsection{Interference-free per-packet solution}
For device $i$, define the isolated service and energy of level $\ell$ as
$b_{i\ell}=\tau B\log_2(1+p_\ell g_i/\sigma^2)$ and
$e_\ell=\tau p_\ell$. Its minimum isolated packet energy is the small integer
problem
\begin{equation}
q_i=\min_{\boldsymbol n\in\mathbb Z_+^{|\mathcal P|}}
 \left\{\sum_\ell n_\ell e_\ell:
 \sum_\ell n_\ell b_{i\ell}\ge L_i,\quad
 \sum_\ell n_\ell\le P_i\right\}.\label{eq:solo}
\end{equation}
Static channels make transmission slots interchangeable for this isolated
problem. We enumerate the counts of all but the lowest positive power and fill
the lowest level minimally. With five levels and periods no larger than 40,
this yields an exact discrete solution. For a fixed evaluated packet set with
$M_i$ packets from device $i$, let $L_0=\sum_i M_iq_i$.

\begin{proposition}\label{prop:bound}
For the zero-loss centralized scheduling problem, $L_0$ is a lower bound on
energy. If $N\le K$, assigning one subchannel to each device attains $L_0$.
\end{proposition}
\begin{proof}
For fixed power, service in \eqref{eq:sinr} cannot exceed isolated service.
Every successfully delivered packet therefore costs at least $q_i$; summing
over packets gives the lower bound. If $N\le K$, distinct subchannels eliminate
interference. Repeating each device's solution of \eqref{eq:solo} within its
nonoverlapping packet windows attains the sum, independently of arrival phases.
\end{proof}

\subsection{Cyclic orthogonal witness at higher loads}
Let $C=\operatorname{lcm}(P_1,\ldots,P_N)$, which is 40 slots in our setup.
For every cyclic packet window $w$ of device $i(w)$, choose integer power-level
counts $n_{w\ell}$. Let $y_{it}\in[0,1]$ indicate slot assignment over one cycle.
An orthogonal schedule is obtained from
\begin{align}
\min_{n,y}\quad &\sum_{w,\ell}e_\ell n_{w\ell}\label{eq:cycle}\\
\text{s.t.}\quad
 &\sum_{t\in w}y_{i(w)t}=\sum_\ell n_{w\ell},\quad \forall w,\nonumber\\
 &\sum_i y_{it}\le K,\quad \forall t,\nonumber\\
 &\sum_\ell b_{i(w)\ell}n_{w\ell}\ge L_{i(w)},\quad\forall w,\nonumber\\
 &n_{w\ell}\in\mathbb Z_+,\quad 0\le y_{it}\le1.\nonumber
\end{align}
For fixed integer counts, the assignment constraints form a capacitated
bipartite flow problem with integral capacities. An integral assignment can
therefore be recovered by maximum flow. Each slot's active devices are assigned
distinct subchannels, and the selected powers are placed within their packet
windows. Repeating this schedule gives a feasible zero-loss witness, whose
energy $U_0$ is measured by replay in the communication environment. We verify
packet completion and energy accounting for every reference instance.

Consequently $L_0\le E_0^\star\le U_0$ for the same zero-loss packet set,
with equality when $N\le K$. Orthogonality is only a restriction for $N>K$;
we make no optimality claim for it there. We report learned packet energy
normalized by both references. Since learned policies may drop packets,
$E/L_0$ and $E/U_0$ characterize their energy relative to zero-loss references,
not a certified approximation factor for the constrained objective
\eqref{eq:objective}. In particular, proximity to the $N=4$ reference does not
establish optimality when a 5\% violation allowance can be used deliberately.
